\documentclass[12pt]{article}
% Basic packages
\usepackage[utf8]{inputenc}   % UTF-8 encoding
\usepackage{geometry}          % Page margins
\usepackage{setspace}          % Double spacing
\usepackage{amsmath, amssymb} % Math symbols
\usepackage{graphicx}         % Include images
\usepackage{booktabs}         % Nice tables
\usepackage{enumitem}         % Custom lists
\usepackage{cancel}
\usepackage{float}

% Page setup: 1-inch margins
\geometry{
  letterpaper,
  margin=1in
}

% Double spacing
\doublespacing{}

% START OF DOCUMENT %
\begin{document}

\begin{center}
    {\textbf{Johnny A. Rojas}}\\
    {\Large \textbf{What Students and Joseph have in common?}}
\end{center}

Good afternoon everyone,

We live in a generation that deeply trusts knowledge.
We study psychology to understand behavior, neuroscience to understand decision-making, and economics to understand success and failure. We are told that with enough research, enough discipline, and enough self-awareness, human beings can learn how to live well.

And in many ways, that pursuit is good.

But something remarkable appears when we look closely.

Modern psychology now tells us that impulsive pleasure weakens long-term satisfaction. That addiction reshapes the brain. That unmanaged stress damages both body and judgment. That delayed gratification predicts success more reliably than intelligence itself.

The question of what truly fulfills human life is not new. Even in ancient Greece, philosophers debated whether happiness came from satisfying desire or learning to restrain it. Thinkers such as Aristippus argued that pleasure itself was the highest good, while Plato warned that a life ruled by impulse ultimately leads not to freedom, but to slavery of the self.

Interestingly, modern psychology has begun revisiting this same question. Research increasingly shows that long-term well-being depends less on satisfying desire and more on the ability to regulate it (Mischel, 1972).

Yet thousands of years before brain scans or behavioral studies existed, Scripture was already warning humanity about these same patterns.

So the question is not whether faith and psychology agree.

The deeper question is this:

**Why do they agree at all?**

If God created human beings — not only spiritually, but biologically and psychologically — then His instructions would not be arbitrary rules. They would function more like operating principles for the human mind itself.

And few stories illustrate this better than the life of Joseph.

Joseph begins his story as many of us do: young, ambitious, full of dreams. But instead of immediate success, his life collapses. He is betrayed by his own brothers, stripped of security, and sold into slavery. (Genesis 37:23–28)

From a psychological perspective, this moment should have defined him.

Research and general knowlege consistently shows that betrayal and loss of control often lead to resentment, anxiety, or destructive coping behaviors. Trauma frequently produces cycles of anger or hopelessness.

Joseph had every justification to become bitter.

And yet, he doesn’t.

Instead of allowing suffering to define his identity, Joseph adapts. He works faithfully even in conditions he did not choose.

Today, psychologists describe this ability to remain stable and purposeful despite hardship as resilience. Developmental psychologist Ann Masten of the University of Minnesota has shown through extensive research that many individuals are capable of adapting positively even after significant adversity, particularly when they maintain responsibility, meaning, and direction in their lives (Masten, 2001).

Joseph’s response follows this same pattern. Rather than allowing betrayal to define him, he continues to work faithfully even in circumstances he did not choose.

But Scripture presents something deeper: Joseph trusted that meaning existed beyond his circumstances.

His stability was not dependent on comfort.

And this becomes even clearer in Potiphar’s house.

Joseph rises to responsibility and influence, but then faces a different kind of trial — temptation. Potiphar’s wife offers him secrecy, pleasure, and immediate reward with seemingly no consequences.

From a modern standpoint, this moment reflects what psychology calls **short-term reward versus long-term identity**.

Contemporary psychological research on self-control, including work by social psychologist Roy Baumeister, shows that human willpower is deeply influenced by environment and repeated exposure to temptation. Individuals are far more successful when they remove themselves from situations that continually test impulse rather than relying on resistance alone (Baumeister et al., 1998).
"

Joseph’s response is fascinating.

He does not negotiate.
He does not remain near temptation.
He runs.

What appears religiously dramatic is psychologically intelligent.

Avoidance of high-risk environments is one of the strongest predictors of behavioral self-control. Recovery programs today teach the same principle: remove proximity to temptation before willpower fails.

Joseph understood something modern research confirms — character is often preserved not by resisting longer, but by leaving sooner.

"After resisting temptation, Joseph is imprisoned unjustly. From every human perspective, integrity appears to have failed him.

Yet modern psychological research suggests that adversity can sometimes produce unexpected development. Psychologists Richard Tedeschi and Lawrence Calhoun describe how individuals frequently experience measurable personal growth following severe hardship, developing stronger identity, resilience, and purpose through suffering itself (Tedeschi & Calhoun, 2004).

What Scripture describes as preparation, psychology now recognizes as growth formed through adversity.

Two languages describing the same reality.

Eventually, Joseph rises to leadership in Egypt — not because suffering disappeared, but because suffering formed capacity.

And here lies the provocative idea:

Perhaps God’s commands are not restrictions meant to limit human freedom.

Perhaps they are descriptions of how humans were designed to function.

When Scripture warns against destructive habits, encourages discipline, promotes rest, generosity, self-control, and faithfulness — it aligns with what psychology increasingly observes about mental health, relational stability, and long-term flourishing.

It is as if creation itself quietly confirms the wisdom behind those commands.

So maybe the conflict we feel is mistaken.

Faith and psychology are not competing explanations of human life.

One studies the design.

The other reveals the Designer.

Joseph’s life suggests that trusting God does not mean abandoning understanding — it means recognizing that the structure of creation, including our own minds, already reflects divine wisdom.

And perhaps the greatest discovery is this:

God does not ask us to live differently from the world simply to test obedience.

He asks us to live according to the way we were made.

Thank you.

\newpage{}

Baumeister, R. F., Bratslavsky, E., Muraven, M., & Tice, D. M. (1998). Ego depletion: Is the active self a limited resource? Journal of Personality and Social Psychology, 74(5), 1252–1265.

Masten, A. S. (2001). Ordinary magic: Resilience processes in development. American Psychologist, 56(3), 227–238.

Mischel, W., Ebbesen, E. B., & Raskoff Zeiss, A. (1972). Cognitive and attentional mechanisms in delay of gratification. Journal of Personality and Social Psychology, 21(2), 204–218.

Tedeschi, R. G., & Calhoun, L. G. (2004). Posttraumatic growth: Conceptual foundations and empirical evidence. Psychological Inquiry, 15(1), 1–18.

\end{document}
